% experiments.tex — Experiments section for the metric-probe study (the "metric angle").
% Meant to be \input into the paper. Needs: amsmath, booktabs. No other packages.
% All numbers are from the completed 500k run matrix and probe battery
% (/scratch/.../factored-fb-metric-probe/), 3 seeds unless noted.

\section{Experiments}\label{sec:experiments}

We test one question: does the bilinear critic succeed because its embedding space encodes
a temporal distance between states? We train the bilinear critic next to three methods that
learn such a distance by construction, probe all of them with the same battery, and check
whether distance quality predicts control performance.

\subsection{Setup}

\paragraph{Environment and data.}
All runs use OGBench \texttt{cube-single-play-v0}: a robot arm must move a cube to a goal
position. The offline dataset holds $10^6$ transitions from 1{,}000 play trajectories of
length 1{,}000. The data is undirected: a scripted policy keeps moving the cube to random
positions and never stops at a goal. Observations are 28-dimensional states and contain no
goal information; actions are 5-dimensional.

\paragraph{Methods.}
We compare five critics. All use the same encoder (3-layer MLP, width 512, LayerNorm), the
same discount $\gamma=0.99$, the same batch size, learning rate, and goal sampling. They
differ only in what the critic learns:
\begin{itemize}
\item \textbf{Bilinear GCIQL} (the subject): $Q(s,a,g)=\langle\varphi(s,a),\psi(g)\rangle/\sqrt d$,
trained with expectile TD learning ($\tau=0.9$); latent size $d\in\{50,512\}$.
\item \textbf{GCIQL}: the same objective with a standard MLP critic on the concatenated
input $(s,a,g)$. This isolates the effect of the bilinear form.
\item \textbf{CRL}: a bilinear critic trained with a contrastive loss against future states
of the same trajectory. Its optimum is the discounted state-occupancy of the behavior
policy.
\item \textbf{QRL}: learns a distance $d_\theta(s,g)$ directly, with an architecture (IQE)
that satisfies the triangle inequality exactly.
\item \textbf{VIP}: learns an embedding $\phi$ and sets $V(s;g)=-\|\phi(s)-\phi(g)\|_2$, so
its distance is symmetric by construction.
\end{itemize}

\paragraph{Policy extraction.}
Every method uses the same flow-matching behavior-cloning policy. For critics that provide
a $Q$ (bilinear GCIQL, GCIQL, CRL) the policy objective adds the critic's value term, as in
one-step actor training. QRL and VIP only learn a representation, so their policy is the
pure behavior-cloning policy. This pure-BC arm doubles as a baseline: it shows how far
cloning the play data gets without any critic.

\paragraph{Controls.}
Two constructions need no training and bound what the probes can show. (i)
\emph{Raw-state distance}: $\|s-g\|_2$ in observation space. Any learned distance must beat
it to claim learned structure, so we also report correlations with this control partialed
out. (ii) \emph{Untrained networks}: the same bilinear architecture at initialization.

\paragraph{Protocol.}
3 seeds per cell, 500{,}000 gradient steps, checkpoints every 100k. Evaluation: OGBench's
five goal-reaching tasks, 50 episodes each; we report mean success over all 250 episodes.
Probes run on a fixed set of arrays drawn once from the held-out validation split (10k
states, 100k state--goal pairs with recorded time gaps, $10^6$ state triples) and every
checkpoint is probed on the identical arrays in float64.

\subsection{Probes}

From each trained critic we build one or more scalar distances $\hat d(s,g)$: for the
bilinear critics, the self-subtracted logit $\hat d(s,g)=f(g,g)-f(s,g)$ and the step count
implied by the value, $\log_\gamma(1+(1-\gamma)V)$; for GCIQL, the step count; for QRL and
VIP, their native distances. We then measure:

\begin{itemize}
\item \textbf{Time correlation.} Spearman rank correlation between $\hat d(s_t,s_{t+\Delta})$
and the time gap $\Delta$ over same-trajectory pairs, with and without partialing out
raw-state distance.
\item \textbf{Triangle inequality.} Over $10^6$ triples $(s,w,g)$: how often
$\hat d(s,g)>\hat d(s,w)+\hat d(w,g)$, and by how much (in steps).
\item \textbf{Symmetry.} $\mathbb{E}|\hat d(s,g)-\hat d(g,s)| \,/\, \mathbb{E}[\hat d(s,g)+\hat d(g,s)]$.
\item \textbf{Dot product vs.\ cosine.} Correlation between the bilinear logit and the
cosine of its two embeddings, plus the spread of embedding norms. If the correlation is
high and norms are tight, the critic is in effect a scaled cosine similarity.
\item \textbf{Effective rank.} Of the embedding matrices over the probe states.
\item \textbf{Cross-method agreement (RSA).} For each method, the matrix of pairwise
distances over a shared 2{,}000-state sample; we report the Spearman correlation between
the upper triangles of every pair of methods. Rank correlation makes this robust to
monotone transforms, so methods that target different quantities can still be compared.
\end{itemize}

\subsection{Results}

\begin{table}[t]
\centering
\caption{Success rate (\%) on \texttt{cube-single} after 500k steps, mean $\pm$ s.d.\ over
3 seeds, and probe summary for the main distance of each method. Time corr.\ is Spearman
$\rho$ with the time gap (in parentheses: partialed on raw-state distance). Triangle
violations report rate / mean size in steps.}
\label{tab:probe}
\small
\begin{tabular}{lccccc}
\toprule
 & Success & Time corr. & Triangle viol. & Symmetry & Eff.\ rank \\
\midrule
Bilinear GCIQL ($d{=}50$) & $\mathbf{92.7\pm1.7}$ & 0.54 (0.41) & 0.9\% / 2.7 & 0.11 & 28 \\
Bilinear GCIQL ($d{=}512$) & $89.2\pm0.7$ & 0.58 (0.46) & 0.3\% / 2.3 & 0.10 & 53 \\
GCIQL (MLP) & $73.2\pm1.6$ & \textbf{0.63 (0.54)} & 0.3\% / 9.5 & 0.06 & --- \\
CRL ($d{=}50$) & $28.1\pm2.1$ & 0.54 (0.45) & 2.2\% / 3.5 & 0.13 & 38 \\
QRL ($d{=}64$) & $8.7\pm1.0$ & 0.55 (0.42) & 6.6\% / ${<}10^{-3}$ & \textbf{0.04} & --- \\
VIP ($d{=}50$) & $9.2\pm2.1$ & 0.04 ($-$0.01) & 0.3\% / 0.1 & 0 & 1.0 \\
\midrule
Raw-state distance & --- & 0.41 (---) & 0 / 0 & 0 & --- \\
Untrained bilinear & --- & 0.02 ($-$0.01) & 48.8\% / 1.5 & $\gg 1$ & 16 \\
Pure BC policy & $8$--$10$ & --- & --- & --- & --- \\
\bottomrule
\end{tabular}
\end{table}

\paragraph{Success.}
Bilinear GCIQL reaches $92.7\pm1.7$, reproducing the historical result. The MLP critic with
the same objective reaches $73.2$. CRL reaches $28$--$34$. The pure-BC policy reaches
$8$--$10$, so on this dataset almost all performance comes from the critic's value term,
and the bilinear form adds 20 points over the MLP critic on top of that. QRL and VIP rows
match the BC baseline by design (no $Q$), so their success numbers only confirm the floor.

\paragraph{The bilinear critic does not have the best distance.}
On every distance probe the MLP critic is as good or better: higher time correlation (0.63
vs.\ 0.54; 0.54 vs.\ 0.41 after removing raw-state distance) and a lower violation rate.
The one probe where the bilinear critic is clearly better is the \emph{size} of triangle
violations: when the MLP critic is wrong it is wrong by 9.5 steps on average, against 2.7
for the bilinear critic. Distance quality measured this way does not track the 20-point
success gap.

\paragraph{Cross-method agreement isolates the bilinear critic.}
The RSA matrix (seed 0) shows that GCIQL, QRL, and CRL agree with each other: GCIQL--QRL
$0.88$, GCIQL--CRL $0.54$, CRL--QRL $0.51$. These three methods, with three different
objectives, recover one shared distance structure. The bilinear critic does not join this
consensus: its agreement with GCIQL is $0.09$, with QRL $0.05$, with CRL $0.21$. Its
nearest neighbor is the raw-state distance ($0.45$--$0.52$). All values are far above the
untrained-network row ($\le 0.04$), so all methods except VIP did learn something --- they
just did not learn the same thing, and the best policy came from the outlier.

\paragraph{The learned solution is a bounded, low-rank, cosine-like critic.}
The bilinear logit correlates with the cosine of its embeddings at $0.996$, and embedding
norms are tight (coefficient of variation $\le 0.09$). We note the untrained network
already shows $0.994$, so this is a property of the LayerNorm architecture, not something
training produces; what training changes is the direction of the embeddings. The effective
rank of the learned embeddings is 28 at $d{=}50$ and 53 at $d{=}512$; success is flat
between the two ($92.7$ vs.\ $89.2$), while $d{=}5$ is known to fail ($9.6$). The solution
the task needs appears to have rank of order 30, and extra capacity above it buys nothing.

\paragraph{Distance structure appears early; success keeps rising.}
Along training of the $d{=}50$ subject (seed 0), the partial time correlation is already
$0.37$ at 100k steps and ends at $0.41$, while success rises from $65.2$ to $93.2$ over the
same interval. Most of the policy improvement happens after the measured distance structure
has stopped changing.

\paragraph{Control outcomes.}
Raw-state distance alone reaches a time correlation of $0.41$ on this environment, which is
why we report partial correlations throughout. The untrained network shows no time
correlation, a $48.8\%$ violation rate, and a symmetry index orders of magnitude above any
trained model. The VIP runs collapsed: effective rank $1.0$ and no time correlation; we
therefore exclude VIP from the comparisons above rather than cite it as a reference point.

\paragraph{Summary.}
If the bilinear critic won by learning a better temporal distance, it should beat the MLP
critic on the distance probes and agree with the methods that provably learn one. It does
neither: it has the weaker distance, it sits outside the consensus the other methods share,
and its geometry stops improving long before its policy does. What it does have is a
bounded, low-rank, cosine-like value function whose errors are small where the MLP critic's
errors are large. This points away from the representation and toward the interface between
critic and policy, which the next section tests directly.
